\chapter{Study B: User Study}
\label{ch:study-b}

% Sources: docs/hipotezler.md (pre-registration), docs/oturum-yonergesi.md (session script),
% data/study/anket_formu.md (questionnaires), data/study/sentences.json (study sentences).
% Written in the past tense before the study was run; every \todo is a value known only afterwards.

Study B asked whether a person can communicate with the full system by gaze alone, and how
abbreviation expansion compares with letter-by-letter gaze typing in the same interface.
It was run with the four members of the project team, so it reports what each participant did
and does not test hypotheses (Section~\ref{sec:b-analysis}).

\section{Expected outcomes}
\label{sec:b-hypotheses}

The pre-registration (Appendix~\ref{app:prereg}) states three expected directions. Because of the
sample size, they are described, not tested.

\begin{description}
  \item[H-B1.] Typing speed (words per minute) is higher with abbreviation expansion (condition B)
        than with letter-by-letter typing (condition A). The technical plan expected an increase of at
        least 30\%; this is reported as each participant's B/A ratio.
  \item[H-B2.] The System Usability Scale score of condition B is above 68.
  \item[H-B3.] Condition B needs fewer corrections per sentence than condition A.
\end{description}

\section{Method}
\label{sec:b-method}

\subsection{Participants}
\label{sec:b-participants}

Four participants took part: the four members of the project team, all undergraduate computer
engineering students (age range \todo{ages}; \todo{number} wore glasses during the session).
The study was first planned with 12--15 volunteers from outside the team; it was changed to the
team before any data were collected (Section~\ref{sec:b-deviations}).

All participants met the pre-registered inclusion criteria: at least 18 years old, native speakers
of Turkish, normal or corrected-to-normal vision, and no known neurological or eye-movement disorder.
Each signed an informed consent form before the session. No camera images were stored; the
recordings contain only numerical features computed from facial landmarks, gaze estimates,
selection events and their times, under a participant code (K01--K04).

Because the participants built the system, they knew the interface, the abbreviation schemes and the
kind of sentences in the corpus. Some had also seen study sentences before their own session; each
participant's exposure was recorded on the experimenter's form (\todo{summary of exposure}).

\subsection{Apparatus}
\label{sec:b-apparatus}

Sessions took place in \todo{room} on a laptop with an NVIDIA RTX~5060 Laptop GPU (8\,GB),
\todo{laptop model, screen size and resolution}. An external 1080p webcam was mounted at the top
centre of the screen, and participants rested their chin on a chin rest at 60\,cm from the screen.
The interface ran full screen in a web browser; face landmarks were detected in the browser with
MediaPipe Face Landmarker, and gaze was mapped to the screen by a calibration model on the server
(Chapter~\ref{ch:system}).

The keyboard occupied the lower 62\% of the screen as a grid of \todo{final grid; 4$\times$3 if kept}
targets. In its first level the grid showed seven letter groups (six of four letters and one of five,
covering the 29 letters of the Turkish alphabet), \emph{Space}, and four actions: \emph{Complete},
\emph{Delete}, \emph{Speak} and \emph{Back}. Selecting a letter took two selections: the group, then
the letter. A target was selected by looking at it for 700\,ms (dwell). The area above the grid showed
the current abbreviation and up to five candidate sentences.

Candidate sentences were generated by Turkish-Gemma-9b fine-tuned with QLoRA (Chapter~\ref{ch:system}),
served on the same laptop by llama.cpp (\todo{model file and version}). Selected sentences were spoken
with the browser's speech synthesis in Turkish (Web Speech API, \texttt{tr-TR}).

\subsection{Design}
\label{sec:b-design}

The study had a within-subjects design with one factor, the input method:

\begin{description}
  \item[A, letter by letter.] The participant typed every letter of the sentence and then selected
        \emph{Speak}. The language model was not used.
  \item[B, abbreviation expansion.] The participant typed only the first letter of each word
        (scheme S1, Chapter~\ref{ch:system}), selected \emph{Complete}, and chose the target sentence from
        the candidates. If it was not among them, the participant could type a word in full and ask
        again.
\end{description}

Everything else was the same in both conditions: the grid, the target size, the dwell time and the
apparatus. Condition order was set by the participant code, not by the experimenter: odd codes did A
then B, even codes B then A, so two participants started with each condition. Codes were drawn by lot.

Each condition was one block of two practice sentences followed by ten measured sentences. There were
two sets of ten measured sentences; the first block always used set~1 and the second set~2, so each
condition met each set in half of the sessions. The task was a copy task: the target sentence was shown
on screen, and a sentence ended only when the spoken sentence matched it.

\subsection{Materials}
\label{sec:b-materials}

The 24 study sentences (20 measured, 4 practice) were short sentences of the kind used in augmentative
and alternative communication, such as requests and questions about care, comfort and daily needs. The
two measured sets were matched in length: 39 and 38 words, 248 and 249 characters (spaces included,
punctuation excluded). None of the study sentences, nor any close variant of them, is in the training,
validation or test part of the corpus; the corpus build checks this automatically.

Rehearsals and the pilot used a separate set of 24 sentences from the training part of the corpus,
matched in length to the study sentences and sharing no word of four or more letters with them, so that
no participant met the study sentences in a rehearsal.

Two questionnaires were used. The Turkish System Usability Scale, SUS-TR
\citep{demirkol2018sustr}, was completed right after the B block; its ten items were used exactly as
published. Five short items written for this study (eye fatigue, unintended selections, effort of
reading the candidates, whether the target sentence was usually among the candidates, and preference
for B over A) were answered once at the end of the session on the same five-point scale. They have no
validation study and are reported item by item.

\subsection{Procedure}
\label{sec:b-procedure}

Each session lasted 40--45 minutes and followed a written script, read aloud word for word
(Appendix~\ref{app:materials}). After consent and a short demographics form, the experimenter explained
both conditions with an example that shares no words with the study sentences. The participant was then
calibrated with nine points, and the calibration was checked on validation points that were not used to
fit it. Two minutes of free practice followed, in which the participant typed one word letter by letter
and deleted a letter; this trained only what the two conditions share, so that getting used to dwell
selection did not fall into the first block.

Then came the first block, a two-minute break, a new calibration, and the second block. SUS-TR was
completed right after the B block, which was the end of the session or the break, depending on the
order. The five short items closed the session.

One author ran all sessions except his own. Since whoever runs a session sees the target sentences on
screen, his own session was run first, by the author who had written the sentences, so that running
sessions did not expose anyone else to the sentences.

During the blocks the experimenter used only fixed phrases. A sentence not finished within three
minutes was skipped. Participants were asked not to talk about the sentences until all sessions were
over.

\subsection{Measures}
\label{sec:b-measures}

Typing speed was measured in words per minute, with the usual convention of five characters per word:
\begin{equation}
  \mathrm{WPM} = \frac{|T| - 1}{t} \times \frac{60}{5},
\end{equation}
where $|T|$ is the number of characters of the target sentence and $t$ is the time in seconds from the
first selection to the speaking of the correct sentence. One is subtracted from $|T|$ because timing
starts with the first selection. Spaces are counted; punctuation is not, because the interface has no
punctuation keys and participants could not enter it.
\todo{Cite the source of this definition after checking it.}

The other measures were: corrections per sentence (selections of \emph{Delete} and \emph{Back}); skipped
sentences; selections per sentence, also expressed as a saving against the fewest selections the same
sentence needs letter by letter in this interface (two per letter, one per space, one for \emph{Speak});
the SUS score; and the five short items.

Calibration error was the mean angle between the gaze estimate and the validation points, which were
not used to fit the calibration. A participant whose error stayed above a threshold after two
recalibrations would have been excluded from both conditions. The threshold was set from the size of the
targets rather than as a fixed number:
\begin{equation}
  \theta = \arctan\!\left(\frac{k/2}{60\,\mathrm{cm}}\right),
\end{equation}
where $k$ is the shorter side of a target on the screen used. Beyond this error the gaze falls, on
average, outside the target being looked at, so the participant cannot use the system.
For the targets used, $k = $ \todo{k} cm and $\theta = $ \todo{$\theta$} degrees.

\subsection{Analysis}
\label{sec:b-analysis}

With four participants, a paired $t$-test ($\alpha = 0.05$, two-sided, power 0.80) detects only effects
of $d_z \approx 2.1$ or larger, which is not a realistic effect size. No hypothesis tests were therefore
run and no $p$-values are reported. For each participant we report the mean WPM in each condition, the
B/A ratio, corrections, skipped sentences and the SUS score, and we summarise them by condition with the
mean and range and by the number of participants who were faster with B.

Skipped sentences are left out of the means. Because skips are more likely in the slower condition,
this makes the slower condition look faster than it was. WPM was therefore also computed with
180\,s, the time limit, for every skipped sentence, and both versions are reported.

All rules in this section were written in the pre-registration before any participant data were
collected.

\section{Results}
\label{sec:b-results}
\todo{Participant flow (calibration, exclusions), per-participant table, summary by condition,
SUS and short items. Results are written after the sessions.}

\section{Changes and deviations from the pre-registration}
\label{sec:b-deviations}

\paragraph{Changes before the pre-registration was locked.}
The first version of the pre-registration (8~November~2026) planned the study with 12--15 volunteers
from outside the team, with a paired test of H-B1 and Holm-corrected tests of H-B2 and H-B3. On
20~November~2026, before any participant data were collected, the team and the advisor decided to run the
study with the four team members, because of the time left before submission. As a consequence, the
hypotheses became expected directions without tests, replacement of excluded participants was dropped,
and the known biases of testing the builders were added to the limitations. In the same period, the
calibration threshold was changed from a fixed 2.5 degrees to the target-size rule above, after the
team's own gaze measurements (1.8--4.5 degrees, one person) showed that a fixed value ignored the target
size. Each change is dated in the pre-registration's change log. The pre-registration was then locked
(\todo{tag onkayit-v1: commit and date}).

\paragraph{Deviations after locking.}
\todo{List every deviation, or state that there were none.}
